\section*{Abstract}
\addcontentsline{toc}{section}{Abstract}

Autonomous agricultural robots operating in dynamic field environments must adapt to
varying terrain conditions to maintain safe and efficient navigation.
This thesis investigates whether terrain type can be reliably inferred from
\emph{intrinsic} sensors alone — accelerometer, gyroscope, and wheel odometry —
without relying on cameras or LiDAR.

A complete machine learning pipeline was designed and implemented, covering raw sensor
quality control, multi-modal synchronisation at 100\,Hz, terrain labelling, signal
cleaning, sliding-window segmentation (2\,s windows, 50\,\% overlap), and the
extraction of over 100 time-domain and frequency-domain features per window.
Six classifier families were evaluated under a leak-safe five-fold group
cross-validation scheme: Logistic Regression, Random Forest, XGBoost, Support Vector
Machine, Multilayer Perceptron (MLP), and a 1D Convolutional Neural Network (CNN).

On a five-class terrain dataset collected across five recording sessions on varied
surfaces (smooth terrain, grass, dry dirt track, muddy dirt track, and cobblestone),
the best model — an MLP operating on the top-25 mutual-information-ranked features —
achieved a macro-averaged F\textsubscript{1} score of 0.766.
Classical models tuned with randomised search reached comparable performance
(Logistic Regression: 0.775 on top-25 features), demonstrating that
expressive feature engineering can largely close the gap between classical and
deep models for this task.
A simplified three-class grouping (smooth, grass, rough) raised the best F\textsubscript{1}
to 0.974, confirming that the principal difficulty lies in distinguishing
visually similar loose-surface classes.

Sensor importance analysis revealed that the gyroscope is the most critical modality:
zeroing gyroscope channels at test time caused a $-$38.7\,\% drop in F\textsubscript{1},
while complete odometry ablation had a negligible impact ($-$0.5\,\%).
Transition detection experiments showed that Random Forest and MLP can identify
terrain changes with a median delay of under one second, even without any explicit
change-point model.

\placeholder{Update abstract with farm data results once collected (expected May–June 2026). Include cross-domain F1, transfer-learning outcome, and key farm-specific findings.}
